\section{Diagnostik Informatif dan Pengujian}
\label{sec:ch16-diagnostik}

\subsection{Kualitas Pesan}

Pesan yang baik memuat:
\begin{itemize}[leftmargin=*]
  \item lokasi yang tepat;
  \item penyebab yang spesifik;
  \item harapan atau aturan yang dilanggar;
  \item saran perbaikan bila memungkinkan.
\end{itemize}

Bandingkan:
\begin{itemize}[leftmargin=*]
  \item buruk: \code{error};
  \item cukup: \code{baris 4: type error};
  \item baik: \code{baris 4, kolom 9: tidak dapat menetapkan float
        ke int; gunakan cast eksplisit}.
\end{itemize}

\subsection{Kode Diagnostik}

Untuk proyek besar, setiap kesalahan dapat memiliki kode, misalnya
\code{E001} untuk identifier tidak dideklarasikan. Kode memudahkan
dokumentasi, pencarian, dan pengujian otomatis.

\subsection{Pengujian}

Pengujian validator mencakup:
\begin{itemize}[leftmargin=*]
  \item program valid harus diterima;
  \item program tidak valid harus ditolak;
  \item pesan kesalahan harus menyebut lokasi yang benar;
  \item jumlah kesalahan harus wajar, tidak berulang.
\end{itemize}

Direktori \file{tests/invalid} memuat program untuk menguji aturan
validasi. Setiap program memiliki satu kesalahan utama, sehingga
hasilnya mudah diverifikasi.

\subsection{Driver Final}

Driver final menggabungkan seluruh fase. Potongannya ditunjukkan pada
\lstref{lst:ch16-main}.

\includecode[
  language=C,
  firstline=40,
  lastline=80,
  firstnumber=40,
  caption={Driver final \file{main.c}},
  label={lst:ch16-main}
]{../../code/subsetc/src/main.c}

Jika tidak ada kesalahan, driver mencetak \code{OK} dan AST. Jika ada
kesalahan, driver mengembalikan status bukan nol. Status ini dapat
digunakan oleh skrip pengujian atau sistem build.

\subsection{M5: Proyek Akhir}

M5 menggabungkan seluruh milestone:
\begin{itemize}[leftmargin=*]
  \item lexer (M1);
  \item parser dan AST (M2);
  \item Bison dan AST lengkap (M3);
  \item symbol table dan type checker (M4);
  \item validator dan diagnostik (M5).
\end{itemize}

Keluaran akhir adalah CLI \code{scc} yang menerima satu berkas
Subset-C, menjalankan seluruh pipeline, dan melaporkan hasilnya.

\begin{table}[H]
\centering
\caption{Artefak M5}
\label{tab:ch16-artefak}
\begin{tabularx}{\textwidth}{@{}l l X@{}}
\toprule
\textbf{Artefak} & \textbf{Status} & \textbf{Keterangan} \\
\midrule
\file{validate.c} & selesai & aturan validasi \\
\file{main.c} & selesai & CLI final \\
\file{tests/invalid} & selesai & tes validasi \\
\file{run-final-tests.bat} & selesai & skrip tes \\
Laporan akhir & wajib & dokumentasi dan refleksi \\
\bottomrule
\end{tabularx}
\end{table}
